\subsection{Constructive algorithm for batched integer-width interchange}
\label{sec:batched-algorithms}

The following pseudocode specifies how the fixed finite circuit becomes
an interchange procedure.  It uses the retained single-digit streaming,
ordered-affine streaming, finite-circuit transfer, and depth-first tape
primitives.

\paragraph{One-time finite preparation in the controlled tensor basis.}
Start from the prime-field five-subset circuit, its corrected role
permutation, and its all-role rational endpoint matrices.  The payload
alphabet is $\mathbb F_3$; the integer-address radix below is a separate
fixed odd prime $q$.

Put $H_0=h^2=784$ and reorder the finite label-space tensor factors so
that the original third factor is first.  Write the resulting space as
$F\otimes V$, where $\dim F=h$ and $\dim V=H_0$, with coordinates
$(\alpha,i)$, $1\le\alpha\le h$, $1\le i\le H_0$, and $i$ varying
fastest.  The controlled basis of
Lemma~\ref{lem:controlled-projector-basis} has the form
\[
 T(u\otimes e_i)=G_i u\otimes e_i,
 \qquad S=T(I_F\otimes K),
\]
where $K\in\operatorname{GL}_{H_0}(\mathbb Q)$ and the
$G_i\in\operatorname{GL}_h(\mathbb Q)$ are chosen independently.
This choice applies to the assigned finite matrices; it introduces no
extra change of coordinates on the tape streams.

\begin{verbatim}
PREPARE_FINITE_DATA(circuit):
    Reorder the finite tensor coordinates as (third, first-and-second).
    Enumerate rational invertible K and rational invertible G_i.
    Form T(u tensor e_i) = G_i*u tensor e_i; S = T*(I tensor K).
    Retain the first choice satisfying the finite controlled-basis tests:
        A1: its first 2h by last 2h corner is nonsingular.
        A3: its first H0+h-1 by last H0+h-1 corner is nonsingular.
        A2: its first H0 by last H0 corner is diagonal and nonsingular.
    Conjugate every terminal and gate matrix by this same S.
    For each edge, form A = matrix_at_head - matrix_at_tail.
    If A belongs to A1 or A3, of rank a:
        r = m-a; t = 2*a-m.
        Store lower-triangular factors A = E1*Pi*E2.
        Store r corner pivots and one middle identity block of t slots.
    If A belongs to A2:
        Store lower-triangular factors A = E1*Pi*E2.
        Store the aligned corner identity block of H0 slots and
            the middle identity block of m-2*H0 slots.
        Store no singleton corner pivots.
    Otherwise:
        Store the ordinary lower-triangular partial-permutation form.
    Choose an odd prime q avoiding all prepared denominators and
        all diagonal numerators needed to keep the factors invertible.
    Form the fixed list of child multipliers t_i in [1, m-1].
    Set tau = 1 - 246/10^9, using the certified strict moment bound.
    Return the finite tables, q, tau, and t_star = 21840.
\end{verbatim}

The corner tests are applied to the conjugated projectors in the three
selected classes.  The controlled-basis lemma proves that the required
rational parameter choice exists, hence that this finite search
terminates.  Every test can be performed by rational elimination.
The lower-triangular factorization normalizes every stored pivot to one.
For the second class,
\[
 A_2=(I_F-\Pi_t)\otimes I_V,
\]
the factor $I_F\otimes K$ commutes with $A_2$.  Its first-to-last block
after conjugation by $T$ is diagonal in $i$, with entries
$[G_i(I_F-\Pi_t)G_i^{-1}]_{1,h}$.  These nonzero entries are normalized
by fixed diagonal units.  They therefore give one identity block in
matching physical order between the first $H_0$ fields of $H$ and the
last $H_0$ fields of $D$.  The Schur complement supplies the second,
contiguous middle identity block.  It is this prescribed matching order
that permits the corner to use a single block interchange.

The common conjugation preserves every edge rank and every all-role
endpoint difference $I_m$.  Ordinary rational elimination constructs
all unselected factorizations.  The prime can be selected by trial
division after the finite forbidden factors are known.  None of this
preparation depends on the input width.

For $m=21952$ and $h=28$, write $B=v^2R$ and $N=v^3$.
The prepared calls for the three selected classes are
\[
\begin{array}{c|c|r|r|l}
 \text{class}&\text{copies}&\text{rank}&\text{singletons}&
                   \text{contiguous block sizes}\\ \hline
 A_1\text{, stage 1--3 join}&B&21896&56&21840\\
 A_2\text{, stage 2 sink}&B&21168&0&784,\ 20384\\
 A_3\text{, stage 3 data entrance}&2N&21141&811&20330
\end{array}
\]
The singleton count, including all ordinary calls on other edges, is
\[
 n_1=s-B(21840+20384+784)-2N\cdot20330
     =105555640096680883200.
\]
Thus the child multipliers are exactly
\[
 1,\quad 784,\quad 21840,\quad 20384,\quad 20330,
\]
with multiplicities $n_1,B,B,B,2N$, respectively.  The exact moment is
\[
 \Psi(x)=\frac{n_1+B(784^x+21840^x+20384^x)+2N\,20330^x}{Wm^x},
 \qquad \Psi(1)=\frac{s}{Wm}.
\]
The controlled-basis certificate proves
\[
 \Psi\!\left(1-\frac{246}{10^9}\right)<1.
\]
All other edges retain the ordinary rank-one implementation, including
any full-rank edge.

\paragraph{Depth and row contract.}
Every recursive invocation below receives an explicit row field preceding
both target chunks.  Its row count $R$ is divisible by $W^{D(e)}$.
Spectator fields of arbitrary lengths may occur elsewhere.

\begin{verbatim}
DEPTH(e):
    depth = 0
    while e >= m:
        e = t_star * floor(e/m)
        depth = depth + 1
    return depth
\end{verbatim}

This computes the function $D$ of Section~\ref{sec:mixed-width-rows}.
It terminates, is nondecreasing, and bounds the number of recursive
circuit invocations on every dependency path.  Every child of a width-$e$
node has depth bound at most $D(e)-1$.

\paragraph{The elementary componentwise shear.}
Suppose two groups each contain $t$ contiguous width-$f$ fields, with the
first group before the second.  The groups can be separated by arbitrary
spectators.  Addition and subtraction in the following routine are
componentwise modulo $q^f$.

\begin{verbatim}
BLOCK_SHEAR(H_block, D_block, t, f, R):
    for i = 1,...,t:
        D_i = D_i + H_i                 # ordered-affine streaming
    INTERCHANGE_WITH_ROWS(H_block, D_block, t*f, R)
    for i = 1,...,t:
        D_i = H_i - D_i                 # ordered-affine streaming
\end{verbatim}

The returned pair is $(H+D,D)$.  The sole recursive call exchanges the
two contiguous blocks without reversing their internal field order.
For $t=1$ the same routine implements an ordinary pivot shear, including
when its two fields have different coordinate indices.

\begin{verbatim}
EDGE_SHEAR(A, H, D, f, R_role):
    Read the stored factors A = E1 * Pi * E2.
    Apply D = E2*D and H = inverse(E1)*H by ordered streaming.
    For each stored singleton pivot (i,j):
        BLOCK_SHEAR(H_i, D_j, 1, f, R_role).
    For each stored identity block (row_start, column_start, t):
        BLOCK_SHEAR(the t consecutive H fields at row_start,
                    the t consecutive D fields at column_start,
                    t, f, R_role).
    Apply H = E1*H and D = inverse(E2)*D by ordered streaming.
\end{verbatim}

This is the address shear $(H,D)\mapsto(H+AD,D)$.
The finite factors are invertible over $\mathbb Z/q^f\mathbb Z$ for
all positive $f$.  Each unbatched pivot uses one width-$f$ call.  Classes $A_1,A_3$
use their singleton corner calls and one middle-block call.  Class $A_2$
uses two block calls, of widths $784f$ and $20384f$, with no singleton
corner calls.  All these widths are strictly less than $mf$.

\paragraph{The finite circuit on complete rows.}
For a main width $mf$, let $H,D$ denote its two $m$-component field
vectors.  All edges below act on one role stream, whose row count is
$R/W$ and logical volume is $V/W$.

\begin{verbatim}
CIRCUIT_SHEAR(H, D, f, R):
    Split each complete row by its index u = W*g + w.
    Process the fixed circuit in chronological order:
        On each edge e:
            EDGE_SHEAR(A_e, H, D, f, R/W) on its role stream.
        At each scalar gate:
            Apply its fixed F_3 table pointwise at matching addresses.
    Merge roles, undoing the known final role permutation.
    Return the original row order.

MAIN_INTERCHANGE(H, D, f, R):
    Apply D = D - H componentwise by ordered-affine streaming.
    CIRCUIT_SHEAR(H, D, f, R)
    Apply D = H - D componentwise by ordered-affine streaming.
\end{verbatim}

The all-role endpoint identity makes \texttt{CIRCUIT\_SHEAR} equal to
$H\leftarrow H+D$ on every original row, including rows assigned to
auxiliary roles.  The three displayed operations in
\texttt{MAIN\_INTERCHANGE} therefore exchange the two main chunks.
All gates allow arbitrary initial auxiliary payloads.  Pointwise
operations preserve the row and spectator coordinates, and the final
merge returns every row to its original index.

\paragraph{An integer width with already available rows.}

\begin{verbatim}
INTERCHANGE_WITH_ROWS(H, D, e, R):
    Require R to be divisible by W^DEPTH(e).
    if e < m:
        Interchange the e corresponding digits individually; return.
    f = floor(e/m);  r = e - m*f.
    Split H = (H_plus,H_minus), D = (D_plus,D_minus),
        with the two high parts of length r.
    Interchange the corresponding high-part digits individually.
    Gather H_plus immediately before H_minus, obtaining
        (D_plus, H_plus, H_minus, gap, D_minus).
    MAIN_INTERCHANGE(H_minus, D_minus, f, R).
    Move H_plus back immediately before the final H_minus.
\end{verbatim}

The result is
$(D^+,D^-,\mathrm{gap},H^+,H^-)$.  Fewer than a fixed multiple of $m$
single-digit operations handle the remainder.  No chunk range is padded
inside this routine.  Every recursive child width is at most
$t_*\lfloor e/m\rfloor$, so splitting the rows once supplies its full
row-divisibility requirement.

\paragraph{An arbitrary radix-$q$ call without a row promise.}

\begin{verbatim}
INTERCHANGE_RADIX_Q(H, D, e):
    depth = DEPTH(e)
    rho = least nonnegative integer with q^(2*rho) >= W^depth
    if e < m or 2*rho >= e:
        Interchange corresponding digits individually; return.
    Split the high rho digits from each chunk.
    Interchange those high parts digit by digit.
    Gather the two high parts before both remaining chunks.
    Regard their joint range q^(2*rho) as an opaque row field.
    Pad complete rows to R = W^depth * ceil(q^(2*rho)/W^depth).
    INTERCHANGE_WITH_ROWS(the remaining chunks, e-rho, R).
    Remove the padded rows.
    Undo the gathering move, preserving the completed high-part swap.
\end{verbatim}

The row range is enlarged by less than two, once.  The existing high
digits supply all its other capacity.  Since $\rho=O(\log e)$, the
fallback condition $2\rho\ge e$ holds for only a fixed finite set of
widths.  Thus its individual-digit cost changes only the constant in
the asymptotic bound.  Computing the displayed powers and comparisons
uses ordinary integer arithmetic; no real logarithm is required at
runtime.  The row count is divisible by $W^{D(e)}$, hence also by
$W^{D(e-\rho)}$.

\paragraph{Binary input chunks.}
For equal binary chunk widths $u$, choose the least integer $e$ with
$q^e\ge2^u$.  Pad each target range numerically from $[2^u]$ to $[q^e]$,
putting zeros on invalid addresses.  Call
\texttt{INTERCHANGE\_RADIX\_Q}, then delete the invalid addresses.
This is the original streaming radix conversion: numerical order is
preserved, each target range expands by less than $q$, and total volume
expands by less than $q^2$.  The completed exchange preserves the set of
valid address pairs.  Since $q$ is fixed and $e=O(u)$, this gives
$O(Vu^\tau)$ for arbitrary binary chunk widths and arbitrary prefix,
gap, and suffix ranges.
